\chapter{Dalsze przykłady ciał i zastosowań}

Ten krótki rozdział zawiera przykłady ilustrujące teorię ciał. Nie stanowi dodatkowego ciągu twierdzeń potrzebnych do głównych celów.

\section{Ciało czteroelementowe i automorfizm Frobeniusa}

Pracujemy tu wyjątkowo w charakterystyce $2$. Wielomian $T^2+T+1\in\F_2[T]$ nie ma zera w $\F_2$, więc jest nierozkładalny. Twierdzenie o ilorazie, ważne w każdej charakterystyce, daje ciało
\[
\F_4=\F_2[T]/(T^2+T+1)=\{0,1,\alpha,\alpha+1\},\qquad \alpha^2=\alpha+1.
\]
Automorfizm Frobeniusa $a\mapsto a^2$ ustala $0,1$ i zamienia $\alpha$ z $\alpha+1$, bo $(\alpha+1)^2=\alpha^2+1=\alpha$. Ma więc rząd $2$. Każdy automorfizm musi wysłać $\alpha$ na jedno z dwóch zer $T^2+T+1$, zatem są dokładnie dwa automorfizmy: identyczność i automorfizm Frobeniusa. Grupa automorfizmów jest więc grupą $C_2$.

\begin{twierdzenie}[Istnienie i jednoznaczność ciał skończonych]
Dla każdej liczby pierwszej $p$ i $m\ge1$ istnieje, z dokładnością do izomorfizmu, dokładnie jedno ciało $\F_{p^m}$ o $p^m$ elementach. Jest ono ciałem rozkładu $T^{p^m}-T$ nad $\F_p$, a jego grupa Galois nad $\F_p$ jest cykliczna rzędu $m$, generowana przez Frobeniusa $\operatorname{Fr}(a)=a^p$.
\end{twierdzenie}
\begin{proof}
W domknięciu algebraicznym $\overline{\F}_p$ rozważmy zbiór $E=\{a:a^{p^m}=a\}$. Pochodna wielomianu $T^{p^m}-T$ wynosi $-1$, więc ma on dokładnie $p^m$ różnych pierwiastków. Ponieważ odwzorowanie $a\mapsto a^{p^m}$ zachowuje sumy i iloczyny, zbiór $E$ jest zamknięty na dodawanie, odejmowanie i mnożenie. Dla $a\ne0$ z $a^{p^m}=a$ wynika $a^{p^m-1}=1$, skąd także $a^{-1}\in E$. Zatem $E$ jest ciałem o $p^m$ elementach.

Jeśli $F$ jest dowolnym ciałem o $p^m$ elementach, jego charakterystyka to $p$ i dla każdego $a\in F^\times$ twierdzenie Lagrange'a daje $a^{p^m-1}=1$. Zatem wszystkie elementy $F$ zerują $T^{p^m}-T$, więc $F$ jest izomorficzne z $E$ (po zanurzeniu w domknięcie algebraiczne $\F_p$).

Frobenius $a\mapsto a^p$ jest automorfizmem $E$ i $\operatorname{Fr}^m=1$. Jego rząd jest dokładnie $m$, gdyż dla $0<d<m$ wielomian $T^{p^d}-T$ ma mniej niż $p^m$ pierwiastków. Natomiast $[E:\F_p]=m$, ponieważ $|E|=p^{[E:\F_p]}$. Liczba automorfizmów nad $\F_p$ nie przekracza stopnia rozszerzenia, a mamy już $m$ różnych potęg Frobeniusa. Są to więc wszystkie automorfizmy.
\end{proof}

\section{Co można odczytać z redukcji wielomianu?}

Dla unormowanego wielomianu $f(T)=\prod_{j=1}^n(T-\alpha_j)\in\Z[T]$ o różnych pierwiastkach definiujemy wyróżnik
$\Delta(f)=\prod_{i<j}(\alpha_i-\alpha_j)^2\in\Z$. Warunek $p\nmid\Delta(f)$ oznacza, że wielomian $\bar f\in\F_p[T]$ jest rozdzielczy.

Aby wykazać związek redukcji z grupą Galois, wykorzystamy krótko liczby $p$-adyczne. Pierścień $\Z_p$ to granica odwrotna pierścieni $\Z/p^k\Z$: element jest zgodnym układem reszt modulo $p,p^2,\ldots$. Mamy naturalną redukcję $\Z_p\to\F_p$, a każde równanie przybliżone z dokładnością do $p^k$ można rozważać dla dowolnie dużego $k$. Pierścień $\Z_p$ jest zupełny w tej topologii, zaś $\Q_p$ oznacza jego ciało ułamków. Użyjemy tylko zupełności i następującego lematu, będącego prostą wersją lematu Hensela.

\begin{lemat}[Proste podnoszenie pierwiastka modulo $p$]
Niech $A$ będzie przemiennym pierścieniem, który jako moduł nad $\Z_p$ jest skończony i wolny, i niech $A/pA$ będzie ciałem. Jeżeli $F\in A[T]$ oraz $\bar b\in A/pA$ spełnia $\bar F(\bar b)=0$ i $\bar F'(\bar b)\ne0$, to istnieje dokładnie jeden $b\in A$ redukujący się do $\bar b$ i spełniający $F(b)=0$.
\end{lemat}
\begin{proof}
Wybieramy podniesienie $b_1$ pierwiastka $\bar b$. Element $F'(b_1)$ jest odwracalny w $A$, ponieważ ma niezerowy obraz w ciele $A/pA$; wystarcza podnieść odwrotność modulo $p$ i użyć zbieżnego szeregu geometrycznego. Definiujemy rekurencyjnie $b_{k+1}=b_k-F(b_k)/F'(b_k)$. Z rozwinięcia Taylora wielomianu wynika, że jeśli $F(b_k)\in p^k A$, to $F(b_{k+1})\in p^{2k}A$. Ponieważ $A$ jest zupełny w topologii $p$-adycznej, ciąg zbiega do zera $F$ o wymaganej redukcji. Jeśli $b,c$ są dwoma takimi zerami, to $F(b)-F(c)=(b-c)H(b,c)$, gdzie $H(b,c)$ ma redukcję $\bar F'(\bar b)\ne0$ i jest odwracalny. Stąd $b=c$.
\end{proof}

\begin{twierdzenie}[Typ cykli z redukcji modulo $p$]
Niech $f\in\Z[T]$ będzie unormowanym wielomianem stopnia $n$ bez pierwiastków wielokrotnych. Jeśli $p\nmid\Delta(f)$ oraz
$\bar f=\bar f_1\cdots\bar f_r$ jest rozkładem na różne unormowane czynniki nierozkładalne w $\F_p[T]$, przy $d_j=\deg\bar f_j$, to $\Gal(f/\Q)$ zawiera element, którego rozłączne cykle na pierwiastkach $f$ mają długości $d_1,\dots,d_r$.
\end{twierdzenie}
\begin{proof}
Pokażemy, skąd pochodzi ten element. Korzystamy z ciała liczb $p$-adycznych $\Q_p$, jego pierścienia $\Z_p$ i poprzedniego lematu; ten krótki argument zastępuje ogólną teorię rozkładu ideałów pierwszych.

Niech $N$ będzie wspólną wielokrotnością $d_1,\dots,d_r$. Ciało $\F_{p^N}$ zawiera wszystkie pierwiastki $\bar f$. Wybieramy element $\bar\theta$ generujący $\F_{p^N}$ nad $\F_p$ i jego wielomian minimalny $\bar h$ stopnia $N$. Podnosimy $\bar h$ do wielomianu unormowanego $h\in\Z_p[T]$ i kładziemy $A=\Z_p[T]/(h)$. Nierozkładalność $\bar h$ zapewnia, że $h$ jest nierozkładalny nad $\Q_p$, a $A/pA\simeq\F_{p^N}$. Ciało ułamków $U=\Q_p[T]/(h)$ ma stopień $N$ nad $\Q_p$.

Wielomian $\bar h$ ma $N$ różnych pierwiastków w $\F_{p^N}$. Poprzedni lemat podnosi każdy z nich jednoznacznie do pierwiastka $h$ w $A$, więc $U/\Q_p$ jest rozszerzeniem Galois. W szczególności podstawienie klasy $\theta$ na pierwiastek $h$ redukujący się do $\bar\theta^p$ definiuje automorfizm $\sigma$ ciała $U$, którego redukcja na $A/pA$ jest Frobeniusem $a\mapsto a^p$.

Wielomian $\bar f$ rozkłada się w $A/pA$ na $n$ różnych czynników liniowych. Znowu stosujemy lemat o podnoszeniu: każdy pierwiastek $\bar f$ ma dokładnie jedno podniesienie będące pierwiastkiem $f$ w $A$. Wszystkie pierwiastki $f$ leżą więc w $U$. Automorfizm $\sigma$ permutuje je, a wskutek jednoznaczności podnoszenia działa na nich tak samo jak Frobenius na pierwiastkach $\bar f$.

Dla nierozkładalnego czynnika $\bar f_j$ stopnia $d_j$ Frobenius działa przechodnio na jego $d_j$ pierwiastkach: jego orbita nad dowolnym pierwiastkiem ma długość równą stopniowi tego pierwiastka nad $\F_p$. Otrzymujemy zatem rozłączne cykle długości $d_1,\ldots,d_r$. Ustalamy zanurzenie globalnego ciała rozkładu $f$ nad $\Q$ do $U$ (wszystkie pierwiastki należą do $U$). Automorfizm $\sigma$ zachowuje to podciało, gdyż permutuje pierwiastki i ustala $\Q$. Jego ograniczenie jest szukanym elementem $\Gal(f/\Q)$.
\end{proof}

\begin{przyklad}[Redukcje wielomianu sześciennego]
Dla $f=T^3-2$ mamy $\Delta(f)=-108$. Modulo $5$ otrzymujemy
$\bar f=(T-3)(T^2+3T+4)$. Drugi czynnik jest nierozkładalny, bo jego wyróżnik wynosi $3$, niebędące kwadratem w $\F_5$. W grupie Galois występuje więc transpozycja. Natomiast modulo $7$ wielomian $T^3-2$ jest nierozkładalny: nie ma zera, gdyż sześciany w $\F_7$ to $0,1,6$. Grupa zawiera zatem również $3$-cykl. To drugi, arytmetyczny dowód, że grupa Galois $T^3-2$ nad $\Q$ jest równa $S_3$.
\end{przyklad}

\begin{uwaga}
W powyższym dowodzie ważny jest warunek $p\nmid\Delta(f)$. Jeśli $\bar f$ ma pierwiastki wielokrotne, jednoznaczne podnoszenie przestaje działać i sam rozkład na czynniki nie daje w ten sposób typu cykli. W dowodzie grupy wielomianu $T^5-4T+2$ nie potrzebowaliśmy tego twierdzenia: wystarczyły nierozkładalność, analiza zer rzeczywistych i elementarny argument o permutacjach.
\end{uwaga}


